\section{Robustness, Mechanisms, and Real Effects}
\label{sec:robustness}

We proceed from the most demanding threat---the 2023 banking
stress---through the specification of the treatment, then turn to where
the deposits go and whether the funding pressure reaches the asset side.

\subsection{The 2023 Stress: Characteristics-by-Quarter Controls}
\label{sec:robustness:stress}

If high-CECL banks were differentially exposed to the March 2023 stress
through observable balance-sheet characteristics, quarter fixed effects
alone would not remove the confound.  Table~\ref{tab:flexcontrols}
re-estimates the baseline uninsured-deposit DiD allowing progressively
richer sets of 2022Q4 bank characteristics to enter interacted with
calendar-quarter fixed effects.  Column~(1) is the baseline.  Column~(2)
adds size and capital interactions; the continuous estimate is essentially
unchanged ($-0.069$, $p<0.10$).  Column~(3) adds the deposit-mix
characteristics most implicated in the 2023 runs (uninsured and brokered
shares), and column~(4) adds the full eight-characteristic vector.  The
level estimates attenuate in columns~(3)--(4) ($-0.043$ and $-0.037$
continuous; $-0.275$ and $-0.238$ binary).

The controls in columns~(3)--(4) include the 2022Q4 uninsured deposit
share, mechanically entangled with the outcome---the bad-control problem
that motivates the composition design.
Table~\ref{tab:flexcontrols_td} applies the same
characteristics-by-quarter saturation \emph{within} the triple
difference; the triple interaction survives essentially intact
($-0.136$ per percentage point of exposure, $p<0.05$; $-1.05$ percentage
points for high-CECL banks, $p<0.01$; baseline: $-0.177$ and $-1.25$).
Whatever differential stress exposure the 2022Q4 characteristics carry,
it does not explain why uninsured deposits fell \emph{relative to
insured deposits within the same bank} in proportion to the disclosed
adjustment.

\subsection{Cohort Timing}
\label{sec:robustness:timing}

The staggered rollout offers a second lever: 323 banks adopted in quarters
other than 2023Q1, experiencing the March 2023 shock at event times other
than $k=0$.  Event studies estimated separately for the early (2022Q3--Q4),
main (2023Q1), and late (2023Q2--Q4) cohorts in event time show the
divergence in each cohort tracking its own adoption quarter rather than
lining up on March 2023, which is difficult to reconcile with a pure
stress story.  The formal test has limited power: dropping the entire
2023Q1 cohort leaves only 13 high-CECL banks, and the resulting DiD
carries the right sign (continuous $-0.081$; binary $-0.133$) but is
imprecise.

A placebo on the 231 late-cohort banks that had not yet adopted as of
2023Q1---regressing the 2023Q1 change in uninsured time deposits on the
size of their still-undisclosed \emph{future} Day-One adjustment---does
not come out clean.  The coefficient is negative and significant
($-0.101$ continuous, $p<0.05$), so late adopters with larger future
adjustments already lost uninsured deposits in the stress quarter.
Two readings are consistent with this pattern.  Under the first, the
March 2023 stress loaded on the same credit-risk fundamentals that CECL
later capitalized into a single number---depositors and CECL agree on
which portfolios are riskier, and the placebo captures this common
fundamental rather than a leakage of undisclosed information.  Under the
second, forward-looking depositors anticipated the disclosure itself in
the weeks between the March failures and the 2023Q1 filings, which would
undermine our characterization of the Day-One adjustment as new
information at adoption.  We do not view the placebo as decisive
between these readings.  What it does establish is that the level
DiD around 2023Q1 cannot by itself separate disclosure from stress at
the primary cohort---which is precisely why identification rests on the
characteristics-by-quarter controls of Section~\ref{sec:robustness:stress}
and on the within-bank composition contrast of the triple difference,
neither of which this cross-sectional 2023Q1 correlation can
contaminate, and on the cohort-timing evidence that response patterns
line up with each cohort's own adoption quarter rather than with March
2023.

\subsection{Extended Pre-Trends, Including the Pandemic Quarters}
\label{sec:robustness:pretrends}

The baseline event window begins in 2020Q3.  Extending the event study
back to 2018Q1 (with and without the COVID-19 quarters), point estimates
over the five-year pre-period oscillate around zero with no economically
meaningful slope, and the post-adoption break remains the dominant
feature.  Joint tests of the pre-adoption coefficients reject marginally
(uninsured $p \approx 0.04$; insured $p \approx 0.06$), driven by
scattered individually significant leads deep in the pre-period rather
than by a trend toward the adoption-date divergence.  The COVID quarters
contribute noise but no differential trend.

\subsection{Residualized Treatment}
\label{sec:robustness:resid}

To isolate variation in the Day-One adjustment orthogonal to everything a
depositor could have observed in 2022, we residualize \texttt{CECL Adj} on
the full expanded predictor set of
Table~\ref{tab:cecl_predictors_cross_section} column~(2) and re-estimate
with the residual as treatment (Table~\ref{tab:resid_cecl}).  Point
estimates retain their signs but attenuate (uninsured deposits $-0.052$
against baseline $-0.072$; interest expense and ROA similarly).  The
depositor response \emph{should} load on the credit-risk content of the
disclosure, and residualization strips out the component of the adjustment
that is correlated with observable credit quality---which is signal, not
confound.  Same sign at roughly two-thirds the magnitude, though
imprecise, suggests the baseline is not driven purely by selection on
observables.

\subsection{Alternative Scalings, Cutoffs, and Nonlinearity}
\label{sec:robustness:scaling}

Scaling the Day-One charge by equity could mechanically link treatment to
capitalization.  Table~\ref{tab:alt_scaling_assets} re-estimates the main
outcomes with the charge scaled by total assets.  The continuous estimates
preserve the baseline pattern: uninsured time deposits decline
($-0.714$, $p<0.10$), interest expense rises, and ROA falls (both
significant at 1\%).  Scaling by loans yields the same pattern
(unreported).  Table~\ref{tab:alt_cutoffs} shows that binary cutoffs at
the 75th and 90th percentiles deliver significant estimates of $-0.28$
and $-0.29$, so the baseline threshold is not doing the work.

The treatment distribution has a mass point at zero, so we examine the
dose-response with economically defined bins.  Figure~\ref{fig:cecl_bins}
plots post-adoption coefficients for uninsured time deposits by adjustment
category, with exactly-zero banks as reference: negative (favorable)
adjustments, small positive, large positive, and extreme positive.  The
response is monotone in the adverse direction: $-0.24$, $-0.40$, and
$-0.56$ across the positive bins.  The left tail is asymmetric: the
negative-adjustment bin also carries a coefficient of $-0.27$
($p<0.05$).  Adverse news drives outflows that scale with its size; any
sizable adjustment in either direction associates with mild outflow
relative to banks reporting exactly zero, consistent with loss-focused
depositor attention.  Dropping the single bank with the largest adjustment
leaves the baseline essentially unchanged (continuous $-0.068$, binary
$-0.381$; unreported).

\subsection{Where Do the Deposits Go?}
\label{sec:robustness:destination}

Are depositors leaving, or re-contracting into insured form?
Table~\ref{tab:brokered} examines the wholesale-funding margin.  After
adoption, high-CECL banks increase insured brokered deposits by 0.52
percentage points of assets ($p<0.05$) and reciprocal deposits---the
market arrangement that purchases full insurance coverage for large
accounts---by 0.089 percentage points per unit of continuous CECL exposure
($p<0.10$), while \emph{uninsured} brokered funding does not move.
Figure~\ref{fig:es_brokered_ins} shows flat pre-trends and post-adoption
accumulation.  High-CECL banks did not lose funding so much as repurchase
it in insured form, at the higher rates documented in
Section~\ref{sec:results:financial}---the market-based insurance margin of
\citet{kim2024insurance}.  Excluding the quartile of banks with the
largest post-adoption growth in insured brokered deposits shrinks the
triple interaction toward zero (unreported), indicating that the
composition shift is concentrated precisely among the banks that actively
replaced uninsured funding with insured wholesale money.

The remaining piece is the local market.  If withdrawn uninsured deposits
re-intermediate nearby, treated-control spillovers would violate the
stable-unit assumption and bias the DiD toward zero.
Table~\ref{tab:spillover} tests this with branch-level SOD data.  Among
low-CECL banks, annual deposit growth after 2023 increases with the
bank's 2022 leave-out exposure to high-CECL competitors in its counties
($+2.4$ percentage points per unit of exposure, $p<0.10$); at the county
level, deposits at low-CECL banks' branches grow faster after 2023 in
counties with larger high-CECL deposit shares ($+5.2$, $p<0.01$).
Deposits leaving high-CECL banks partly land at local competitors, the
direction of spillover that makes our within-market estimates
conservative.

\subsection{Lending and the Bank-Response Channel}
\label{sec:robustness:lending}

Banks respond to CECL through provisioning, capital management, and
credit supply \citep{granja2025current, chen2022does, kim2026current}.
Could deposit movements reflect the bank's actions rather than
depositors' reading of the disclosure?  Three observations bound this
channel: deposit outflows appear simultaneous with the first possible
strategic response and before an altered lending policy could plausibly
be observed by depositors; flat pre-trends rule out the version in which
banks adjusted ahead of adoption; and the funding pressure
\emph{precedes} the lending response.
Table~\ref{tab:lending} shows quarterly total loan growth at high-CECL
banks slows by 0.65 percentage points ($p<0.01$), concentrated in
commercial real estate ($-0.90$, $p<0.01$).
Figure~\ref{fig:es_loan_growth} shows the slowdown building over the
post-period rather than jumping at adoption, as expected if lending
adjusts to a funding shock with a lag---the funding-cost pass-through
channel emphasized by \citet{granja2025current}.
